\subsubsection{Herramientas para el desarrollo del modelo}

El modelo de clasificación del sistema se desarrollará utilizando
Python como lenguaje principal, apoyándose en un conjunto de
herramientas especializadas en aprendizaje profundo y visión por
computadora.

\textbf{TensorFlow y Keras} constituirán el núcleo del proceso de
entrenamiento. TensorFlow es un framework de código abierto desarrollado
por Google, ampliamente utilizado en la industria y la academia para
la construcción y entrenamiento de modelos de aprendizaje profundo.
Keras actúa como interfaz de alto nivel sobre TensorFlow, simplificando
la definición de arquitecturas de redes neuronales convolucionales como
EfficientNet-B3, seleccionada para la clasificación de imágenes de
árboles en este proyecto.

\textbf{TensorFlow Lite} es la versión optimizada de TensorFlow
orientada a dispositivos móviles y embebidos. Una vez entrenado el
modelo con Keras, este se convertirá al formato \texttt{.tflite},
reduciendo su tamaño y optimizando su rendimiento para ejecutarse
directamente en el dispositivo Android sin requerir conexión a internet.
Esta conversión es fundamental para el enfoque offline-first del sistema.

\textbf{Scikit-learn} se empleará como herramienta complementaria
durante el proceso de entrenamiento y evaluación del modelo. Aportará
métricas de evaluación como precisión, recall, F1-score y matriz de
confusión, permitiendo analizar el desempeño del modelo de clasificación
de forma rigurosa antes de su conversión a formato TFLite.

\textbf{Amazon SageMaker} se empleará como plataforma de registro y
versionado del modelo. Una vez que el modelo sea entrenado en los
equipos de cómputo disponibles, se registrará en SageMaker Model
Registry para llevar control de versiones y garantizar la trazabilidad
de cada iteración antes de publicarlo en Amazon S3 para su distribución
al dispositivo móvil.

El entrenamiento se realizará sobre los equipos de cómputo del equipo
de desarrollo, aprovechando las capacidades de procesamiento local
disponibles para la experimentación y ajuste de hiperparámetros del
modelo.